\chapter{激光稳频}
\section{饱和吸收稳频}
\section{调制转移稳频}
调制转移谱（Modulation Transfer Spectroscopy, MTS）是一种广泛应用于激光稳频的高灵敏光谱技术。其核心思想是通过电光相位调制在激光光场上产生边带，再利用原子的饱和吸收效应将调制信息转移到探测光的幅度上，最终通过解调获得色散型误差信号，实现激光频率对原子跃迁线的锁定。

首先，使用电光调制器（EOM）对频率为 $\omega_0$ 的单频激光进行相位调制。调制后的光场可表示为：
\begin{equation}
	E(t) = E_0 e^{i\left(\omega_0 t + \beta \sin(\Omega_m t)\right)},
\end{equation}
其中 $E_0$ 为光场振幅，$\beta$ 为调制深度，$\Omega_m$ 为调制频率。利用贝塞尔函数恒等式 $e^{i z \sin\theta} = \sum_{n=-\infty}^{\infty} J_n(z) e^{i n \theta}$，可将光场展开为：
\begin{equation}
	E(t) = E_0 \sum_{n=-\infty}^{\infty} J_n(\beta) e^{i\left(\omega_0 + n \Omega_m\right) t},
\end{equation}
其中 $J_n(\beta)$ 为第 $n$ 阶贝塞尔函数。上式表明，相位调制后的激光包含了载波 $\omega_0$ 以及一系列边带 $\omega_0 \pm n \Omega_m$。


将调制后的激光分为两束：**泵浦光**（较强）和**探测光**（较弱），两束光反向传播通过原子/分子样品池。
当激光频率扫描至原子跃迁频率 $\omega_a$ 附近时，强泵浦光会使原子激发态饱和，在多普勒展宽的吸收谱线上“烧”出一个窄孔（饱和吸收效应）。
由于泵浦光带有相位调制边带，其与探测光在原子介质中的非线性相互作用（四波混频过程）会将泵浦光的调制信息转移到探测光的幅度上。此时，探测光不仅经历吸收，其幅度也会携带频率为 $\Omega_m$ 的调制成分。

通过光电探测器探测经过样品池后的探测光，并将信号输入锁相放大器（Lock-in Amplifier），以调制频率 $\Omega_m$ 为参考进行解调。
解调后得到的信号是一个**色散型（Dispersive）误差信号**：
- 当激光频率 $\omega_0 = \omega_a$ 时，误差信号为零；
- 当 $\omega_0 > \omega_a$ 时，误差信号为正；
- 当 $\omega_0 < \omega_a$ 时，误差信号为负。

将此误差信号反馈回激光的频率调谐机构（如压电陶瓷、温度控制器），即可将激光频率稳定在原子跃迁线的中心频率 $\omega_a$ 处。

调制转移谱稳频具有灵敏度高、谱线窄、无多普勒背景等优点，是冷原子物理、量子光学等领域中常用的激光稳频技术。
\section{PDH稳频}